% Part: many-valued-logic
% Chapter: three-valued-logics
% Section: lukasiewicz

\documentclass[../../../include/open-logic-section]{subfiles}

\begin{document}

\olfileid{mvl}{thr}{luk}

\olsection{د لوکاشېوېچ منطق}

د څو ارزښته منطق له لومړنيو خپرو شوو او بشپړ تشريح شوو وړانديزونو
څخه يو د پولنډي فيلسوف يان لوکاشېوېچ و، چې په ۱۹۲۱ کال کښې ئې
وړاندې کړ۔ د ارسطو د سمندري جګړې مسئلې هغه دې ته وهڅاوه: داسې
ښکاري چې \emph{نن} د «سبا به سمندري جګړه وي» جمله نه رښتيا ده
او نه دروغ؛ د هغې د رښتياوالي قيمت لا نۀ دے ټاکل شوے۔ لوکاشېوېچ
وړانديز وکړ چې د دا ډول «راتلونکو اتفاقي» جملو دپاره درېيم
د رښتياوالي قيمت ورزيات شي۔
\begin{quote}
زه له تناقض پرته دا فرض کولے شم چې راتلونکے کال په يوه ټاکلې
شېبه کښې، مثلاً د دسمبر پر ۲۱مه غرمه، په وارسا کښې زما شتون
اوس نه په مثبت ډول ټاکل شوے دے او نه په منفي ډول۔ نو په هغه
وخت کښې زما په وارسا کښې شتون ممکن دے، خو لازم نۀ دے۔ د دې
فرض له مخې «زه به راتلونکے کال د دسمبر پر ۲۱مه غرمه په وارسا
کښې يم» بيان اوس نه رښتيا کېدے شي او نه دروغ۔ که دا اوس رښتيا
وے، نو په راتلونکي کښې زما په وارسا کښې شتون به لازم وے، او
دا له فرض سره په تناقض کښې دے۔ له بلې خوا، که دا اوس دروغ وے،
نو په راتلونکي کښې زما هلته شتون به ناشونے وے، او دا هم له فرض
سره په تناقض کښې دے۔ نو تر غور لاندې بيان اوس نه رښتيا دے او
نه دروغ؛ درېيم قيمت غواړي چې له «۰» يا دروغو او «۱» يا رښتيا
څخه بېل وي۔ دا قيمت په «$\frac{1}{2}$» ښودلے شو۔ دا «ممکن»
ښيي او د «رښتيا» او «دروغ» تر څنګ درېيم قيمت جوړوي۔
\end{quote}
موږ به د لوکاشېوېچ د درېيم رښتياوالي قيمت دپاره~$\Undef$
کاروو۔\footnote{لوکاشېوېچ دلته «ممکن» په هغه معنا کاروي چې نن
عامه نۀ ده: ممکن، خو لازم نۀ۔}

په دې تفسير کښې د~$\lnot$، $\land$ او~$\lor$ نښلوونکو
د صدق تابعې ټاکل اسانه دي: د راتلونکي اتفاقي بيان نفي هم
راتلونکے اتفاقي بيان دے، نو~$\tf{\lnot}(\Undef) = \Undef$۔
که د عطف يو اړخ ناټاکلے او بل رښتيا وي، ټول عطف هم ناټاکلے
وي؛ ځکه چې د راتلونکي اتفاقي اړخ له پايلې سره سم عطف کېدے شي
رښتيا شي او کېدے شي دروغ شي۔ نو \[
    \tf{\land}(\True, \Undef) =
\tf{\land}(\Undef, \True) =
\Undef.
\]
خو که بل اړخ دروغ وي، عطف رښتيا کېدے نۀ شي؛ نو
\[\tf{\land}(\False, \Undef) =
\tf{\land}(\Undef, \False) = \False.\]
د سرچينې يادونه: د دويمې جوړې د مدخلونو ترتيب په اصلي ښودنه کښې
په تېروتنه تکرار شوے و؛ دلته دواړه ترتيبونه ښودل شوي دي۔
نور قيمتونه، چې مدخلونه ئې ټاکلي رښتياوالي قيمتونه~$\True$
يا~$\False$ وي، د کلاسيک منطق په څېر دي۔

د شرطي په باب حالت لږ پېچلے دے۔ فرض کړئ~$q$ يو راتلونکے
اتفاقي بيان دے۔ که~$p$ دروغ وي، نو~$p \lif q$ به رښتيا وي،
د~$q$ له راتلونکې پايلې پرته؛ نو بايد
$\tf{\lif}(\False, \Undef) = \True$ وټاکو۔ او که~$p$
رښتيا وي، نو~$q \lif p$ به د~$q$ له پايلې پرته رښتيا وي؛
نو~$\tf{\lif}(\Undef, \True) = \True$۔ که~$p$ رښتيا
وي، نو~$p \lif q$ کېدے شي رښتيا يا دروغ شي؛ نو
$\tf{\lif}(\True, \Undef) = \Undef$۔ په همدې ډول، که~$p$
دروغ وي، نو~$q \lif p$ کېدے شي رښتيا يا دروغ شي؛ نو
$\tf{\lif}(\Undef, \False) = \Undef$۔ هغه حالت پاتې کېږي
چې~$p$ او~$q$ دواړه راتلونکي اتفاقي بيانونه وي۔ د دې انګېزې
له مخې، په دې حالت کښې بايد رښتيا~$\Undef$ ورکړو۔ خو په دې سره
$!A \lif !A$ به \emph{نۀ} تاوتولوژي وي۔ لوکاشېوېچ د
$!A \lor \lnot !A$ او~$\lnot(!A \land \lnot !A)$ له پرېښودو
سره ستونزه نۀ لرله، خو د~$!A \lif !A$ پرېښودل ئې نۀ منل۔
نو ئې~$\tf{\lif}(\Undef, \Undef) = \True$ وټاکه۔

\begin{defn}\ollabel{def:lukasiewicz}
درې ارزښته لوکاشېوېچ منطق د لاندې ماتريس په وسيله تعريفېږي:
\begin{enumerate}
  \item معياري بياني ژبه~$\Lang L_0$ چې
  $\lnot$, $\land$, $\lor$, $\lif$
  نښلوونکي لري۔
  \item د رښتياوالي د قيمتونو سټ~$V = \{\True, \Undef, \False\}$۔
  \item $\True$ يوازينے ټاکل شوے قيمت دے؛ يعنې~$V^+ = \{\True\}$۔
  \item د صدق تابعې په لاندې جدولونو کښې ورکړل شوې دي:
  \begin{center}
    \begin{tabular}{c|c} 
      $\tf{\lnot}$ & \\ 
      \hline  
      $\True$ & $\False$ \\ 
      $\Undef$ & $\Undef$ \\
      $\False$ & $\True$ 
    \end{tabular}
    \quad
    \begin{tabular}{c|ccc} 
      $\tf{\land}[\LogLuk[3]]$ & $\True$ & $\Undef$ & $\False$ \\ 
      \hline 
      $\True$ & $\True$ & $\Undef$ & $\False$ \\ 
      $\Undef$ & $\Undef$ & $\Undef$ & $\False$\\ 
      $\False$ & $\False$ & $\False$ & $\False$ 
    \end{tabular}
    \\[2ex]
    \begin{tabular}{c|ccc} 
      $\tf{\lor}[\LogLuk[3]]$ & $\True$ & $\Undef$ & $\False$ \\ 
      \hline 
      $\True$ & $\True$ & $\True$ & $\True$ \\ 
      $\Undef$ & $\True$ & $\Undef$ & $\Undef$ \\
      $\False$ & $\True$ & $\Undef$ & $\False$ 
    \end{tabular}
    \quad
    \begin{tabular}{c|ccc} 
      $\tf{\lif}[\LogLuk[3]]$ & $\True$ & $\Undef$ & $\False$ \\ 
      \hline 
      $\True$ & $\True$ & $\Undef$ & $\False$ \\ 
      $\Undef$ & $\True$ & $\True$ & $\Undef$  \\ 
      $\False$ & $\True$ & $\True$ & $\True$ 
    \end{tabular}
  \end{center} 
\end{enumerate}
\end{defn}

لکه په اسانه چې څرګندېږي، هر داسې !!{formula}~$!A$ چې يوازې
$\lnot$، $\land$ او~$\lor$ لري، هغه وخت د رښتياوالي
قيمت~$\Undef$ اخلي چې ټولو !!{propositional variable}s ته ئې
$\Undef$ ورکړل شوے وي۔ نو، د بېلګې په توګه، کلاسيکې
تاوتولوژۍ~$p \lor \lnot p$ او~$\lnot(p \land \lnot p)$
په~$\LogLuk[3]$ کښې تاوتولوژۍ نۀ دي، ځکه چې کله
$\pAssign v(p) = \Undef$ وي، نو~$\pValue{v}(!A) = \Undef$۔

په هغو !!{valuation}s کښې چې د فارمول د هر بياني متغير دپاره
$\pAssign v(p) = \True$ يا~$\False$ وي، $\pValue v(!A)$
له کلاسيک رښتياوالي قيمت سره برابر وي۔

\begin{prop}
  که د~$!A$ د هر~$p$ دپاره
  $\pAssign v(p) \in \{\True, \False\}$ وي، نو
  $\pValue v(!A)[\LogLuk[3]] = \pValue v(!A)[\LogCL]$.
\end{prop}

\begin{prob}\label{mvl:thr:luk:prob:luk-iff} فرض کړئ
  $\pValue v(!A \liff !B) = \pValue v((!A \lif !B) \land (!B \lif
  !A))$ په~$\LogLuk[3]$ کښې داسې تعريفوو۔ د~$\liff$
  د رښتياوالي جدول به څنګه وي؟
\end{prob}

ډېرې کلاسيکې تاوتولوژۍ په~\LogLuk[3] کښې هم \emph{تاوتولوژۍ دي}،
لکه~$\lnot p \lif (p \lif q)$۔ د کلاسيک منطق په څېر، دا د
رښتياوالي په جدولونو هم تصديق کولے شو:
\begin{center}
\begin{tabular}{cc|cccccc}
  $p$ & $q$ & $\lnot$ & $p$ & $\lif$ & $\smash{(}p$ & $\lif$ & $q\smash{)}$ \\ \hline
  \True & \True & \False & \True & \True & \True & \True & \True \\
  \True & \Undef & \False & \True & \True & \True & \Undef & \Undef \\
  \True & \False & \False & \True & \True & \True & \False & \False \\
  \Undef & \True & \Undef & \Undef & \True & \Undef & \True & \True \\
  \Undef & \Undef & \Undef & \Undef & \True & \Undef & \True & \Undef \\
  \Undef & \False & \Undef & \Undef & \True & \Undef & \Undef & \False \\
  \False & \True & \True & \False & \True & \False & \True & \True \\
  \False & \Undef & \True & \False & \True & \False & \True & \Undef \\
  \False & \False & \True & \False & \True & \False & \True & \False \\  
\end{tabular}
\end{center}

\begin{prob}
  وښيئ چې لاندې فارمولونه په~\LogLuk[3] کښې تاوتولوژۍ دي:
  \begin{enumerate}
    \item $p \lif (q \lif p)$
    \item\label{mvl:thr:luk:prob:luk-taut-2} $\lnot(p \land q) \liff (\lnot p \lor \lnot q)$
    \item\label{mvl:thr:luk:prob:luk-taut-3} $\lnot(p \lor q) \liff (\lnot p \land \lnot q)$
  \end{enumerate}
  (په~\olref[mvl][thr][luk]{prob:luk-taut-2} او
  \olref[mvl][thr][luk]{prob:luk-taut-3} کښې $!A \liff !B$
  د~$(!A \lif !B) \land (!B \lif !A)$ د لنډيز په توګه واخلئ،
  يا د~\olref[mvl][thr][luk]{prob:luk-iff} خپل حل ته رجوع وکړئ۔)
\end{prob}

\begin{prob}
  وښيئ چې لاندې کلاسيکې تاوتولوژۍ په~\LogLuk[3] کښې
  تاوتولوژۍ نۀ دي:
  \begin{enumerate}
    \item $(\lnot p \land p) \lif q$
    \item $((p \lif q) \lif p) \lif p$
    \item $(p \lif (p \lif q)) \lif (p \lif q)$
  \end{enumerate}
  د سرچينې يادونه: په لومړي فارمول کښې يو اضافي تړونيز قوس
  و؛ دلته هغه لرې شوے دے۔
\end{prob}

له دې څخه ښايي څوک داسې وانګېري چې که څه هم ټولې کلاسيکې
تاوتولوژۍ په~$\LogLuk[3]$ کښې تاوتولوژۍ نۀ دي، لږ تر لږه بايد
په هره !!{valuation} کښې يا~$\True$ يا~$\Undef$ قيمت واخلي۔
خو داسې نۀ ده۔ لاندې فارمول ضد مثال دے:
\[
  \lnot(p \lif \lnot p) \lor \lnot(\lnot p \lif p)
\]
که~$p$، $\Undef$ وي، د دې قيمت~$\False$ وي۔

\begin{prob}
  لاندې کومې اړيکې د لوکاشېوېچ په منطق کښې صدق کوي؟ د هرې
  يوې دپاره د رښتياوالي جدول ورکړئ۔
  \begin{enumerate}
    \item $p, p \lif q \Entails q$
    \item $\lnot\lnot p \Entails p$
    \item $p \land q \Entails p$
    \item $p \Entails p \land p$
    \item $p \Entails p \lor q$
  \end{enumerate}
\end{prob}

لوکاشېوېچ هيله لرله چې د خپل درې ارزښته نظام پر بنسټ د امکان
منطق جوړ کړي۔ د دې دپاره ئې يوځايه نښلوونکے~$\Diamond !A$
(د «$!A$ ممکن دے» دپاره) او ورسره اړوند~$\Box !A$
(د «$!A$ لازم دے» دپاره) ورزيات کړل:
\begin{center}
  \begin{tabular}{c|c} 
    $\tf{\Diamond}$ & \\ 
    \hline  
    $\True$ & $\True$ \\ 
    $\Undef$ & $\True$ \\
    $\False$ & $\False$ 
  \end{tabular}
  \quad
  \begin{tabular}{c|c} 
    $\tf{\Box}$ & \\ 
    \hline  
    $\True$ & $\True$ \\ 
    $\Undef$ & $\False$ \\
    $\False$ & $\False$ 
  \end{tabular}
\end{center}
په بل عبارت، $p$ هله او يوازې هله ممکن دے چې لا ئې دروغوالے
نۀ وي ټاکل شوے؛ او~$p$ هله او يوازې هله لازم دے چې رښتياوالے
ئې لا ټاکل شوے وي۔

\begin{prob}
  وښيئ چې~$\Box p \liff \lnot\Diamond \lnot p$ او
  $\Diamond p \liff \lnot \Box \lnot p$ په هغه~\LogLuk[3]
  کښې تاوتولوژۍ دي چې د~$\Box$ او~$\Diamond$ د
  رښتياوالي جدولونه ورزيات شوي وي۔
\end{prob}

خو د دې وړانديز شوي مودالي منطق نيمګړتياوې ژر څرګندې شوې:
هر څه چې وشي، $p \land \lnot p$ هېڅکله رښتيا کېدے نۀ شي۔
نو که څه هم اوس ئې ارزښت نۀ وي ټاکل شوے (او ځکه ناټاکلے وي)،
بايد ناشونے وګڼل شي؛ يعنې~$\lnot \Diamond(p \land \lnot p)$
بايد تاوتولوژي وي۔ خو که~$\pAssign v(p) = \Undef$ وي، نو
$\pValue v(\lnot \Diamond(p \land \lnot p)) = \False$۔
د سرچينې يادونه: اصلي جمله دلته د نتيجې قيمت ناټاکلے
ليکي، خو د پورته صدق جدولونو له مخې دا قيمت دروغ دے؛
په دواړو صورتونو کښې فارمول تاوتولوژي نۀ ده۔ که څه هم
لوکاشېوېچ په دې کښې سم و چې د امکان او لزوم په شان مودالي
توپيرونو دپاره دوه د رښتياوالي قيمتونه بس نۀ دي، د درېيم
قيمت ورزياتول هم بس نۀ دي۔

\end{document}
